%! Polynomial Functions and End Behavior — Unit 2, Lesson 2.5 — Slides (Beamer)
\documentclass[aspectratio=169, 11pt]{beamer}
\usepackage{precalculus-beamer}   % provides \plumheader{} and \sectionlabel[color]{}

\begin{document}

% TITLE SLIDE
{
\setbeamercolor{background canvas}{bg=plum}
\begin{frame}[plain]
  \vspace{2.8cm}\hspace{0.55cm}%
  \begin{minipage}{0.9\textwidth}
    {\color{goldacc}\bfseries\small LESSON 2.5}\\[0.5cm]
    {\color{white}\bfseries\LARGE Polynomial Functions and End Behavior}\\[1.2cm]
    {\color{lilacmid}\small Precalculus~$\cdot$~Unit 2: Polynomial Functions}
  \end{minipage}
\end{frame}
}

% HOOK SLIDE
\begin{frame}[t, plain]
  \plumheader{Hook: The Graph That Lies}
  \sectionlabel[goldacc]{Think About It}

  \begin{columns}[T]
    \column{0.5\textwidth}
      \begin{center}
      \begin{tikzpicture}[x=1.0cm, y=1.1cm]
        \draw[linegray, very thin, xstep=0.5, ystep=0.5] (-1,-1.5) grid (3.5,0.5);
        \draw[->, thick] (-1.3,0) -- (3.9,0) node[below right] {\small \(x\)};
        \draw[->, thick] (0,-1.8) -- (0,0.9) node[above] {\small \(p(x)\)};
        \foreach \x/\lab in {1/20, 2/40, 3/60} \node[above] at (\x,0.08) {\small \lab};
        \draw[very thick, plum] plot [smooth] coordinates
          {(-1,-0.48) (-0.5,-0.11) (0,0) (0.5,-0.09) (1,-0.32) (1.5,-0.63)
           (2,-0.96) (2.5,-1.25) (3,-1.44) (3.33,-1.48) (3.5,-1.47)};
        \fill[plum] (0,0) circle (2.5pt);
      \end{tikzpicture}
      \end{center}

    \column{0.46\textwidth}
      \(p(x) = x^3 - 100x^2\) on \(-20 \le x \le 70\).

      \bigskip
      \begin{itemize}\setlength{\itemsep}{8pt}
        \item It dives to about \(-150{,}000\) and flattens out.
        \item It looks like a function that has given up.
        \item Predict: is \(p(1000)\) big and negative?
      \end{itemize}

      \bigskip
      \textit{\(p(1000) = 900{,}000{,}000\). Positive --- nearly a billion.
      Nothing was wrong with the window. It was just too small.}
  \end{columns}
\end{frame}

% SECTION 1 — LIMIT NOTATION
\begin{frame}[t, plain]
  \plumheader{1. Limit Notation --- Language for the Two Far Ends}

  \begin{block}{The behavior, in words}
    As the input values of a nonconstant polynomial increase without bound, the
    output values either \textbf{increase without bound} or \textbf{decrease
    without bound}. There is no third option --- a polynomial cannot level off.
  \end{block}

  \medskip
  \begin{columns}[T]
    \column{0.48\textwidth}
      \begin{block}{Running right}
        \[ \lim_{x \to \infty} p(x) = \infty
           \qquad \lim_{x \to \infty} p(x) = -\infty \]
      \end{block}

    \column{0.48\textwidth}
      \begin{block}{Running left}
        \[ \lim_{x \to -\infty} p(x) = \infty
           \qquad \lim_{x \to -\infty} p(x) = -\infty \]
      \end{block}
  \end{columns}

  \bigskip
  \textit{\(\infty\) is not a number. You never substitute it and never get it
  back --- the statement \textbf{describes} a behavior, it does not compute one.
  Never write \(p(\infty) = \infty\).}
\end{frame}

% SECTION 2 — DOMINANCE
\begin{frame}[t, plain]
  \plumheader{2. Why the Leading Term Wins}

  \begin{block}{The claim --- and it is a \textit{because}}
    Degree and the sign of the leading term decide end behavior \textbf{because}
    for inputs of large magnitude, the leading term dominates the values of
    every lower-degree term.
  \end{block}

  \medskip
  \begin{center}
  \begin{tabular}{|r|c|c|c|c|}
  \hline
  \(x\) & \textbf{leading: \(x^3\)} & \textbf{the rest: \(-100x^2\)} & \(p(x)\) & \textbf{rest \(\div\) leading} \\
  \hline
  \(10\)      & \(1 \times 10^{3}\)    & \(-1 \times 10^{4}\)   & \(-9 \times 10^{3}\)    & \(-10\) \\
  \hline
  \(50\)      & \(1.25 \times 10^{5}\) & \(-2.5 \times 10^{5}\) & \(-1.25 \times 10^{5}\) & \(-2\) \\
  \hline
  \(100\)     & \(1 \times 10^{6}\)    & \(-1 \times 10^{6}\)   & \(0\)                   & \(-1\) \\
  \hline
  \(500\)     & \(1.25 \times 10^{8}\) & \(-2.5 \times 10^{7}\) & \(1 \times 10^{8}\)     & \(-0.2\) \\
  \hline
  \(1{,}000\) & \(1 \times 10^{9}\)    & \(-1 \times 10^{8}\)   & \(9 \times 10^{8}\)     & \(-0.1\) \\
  \hline
  \end{tabular}
  \end{center}

  \bigskip
  \textit{The last column is \(\dfrac{-100x^2}{x^3} = \dfrac{-100}{x} \to 0\).
  The rest never gets small --- it gets small \textbf{compared with} the leading
  term.}
\end{frame}

% SECTION 3 — THE FOUR CASES
\begin{frame}[t, plain]
  \plumheader{3. The Four Cases}

  \begin{block}{Only two questions matter}
    Is the degree \textbf{even or odd}? Is the leading coefficient
    \textbf{positive or negative}? Two choices twice over --- four cases, and
    there are no others.
  \end{block}

  \bigskip
  \begin{center}
  \begin{tabular}{|c|c|c|c|c|}
  \hline
   & \textbf{Degree} & \textbf{Leading coeff.} & \(\lim_{x \to \infty} p(x)\) & \(\lim_{x \to -\infty} p(x)\) \\
  \hline
  Case 1 & Even & Positive & \(\infty\)  & \(\infty\) \\
  \hline
  Case 2 & Even & Negative & \(-\infty\) & \(-\infty\) \\
  \hline
  Case 3 & Odd  & Positive & \(\infty\)  & \(-\infty\) \\
  \hline
  Case 4 & Odd  & Negative & \(-\infty\) & \(\infty\) \\
  \hline
  \end{tabular}
  \end{center}

  \bigskip
  \textit{Even degree sends both ends the same way; odd degree sends them
  opposite ways; a negative leading coefficient flips both.}
\end{frame}

% SECTION 3, CONT. — THE FOUR GRAPHS
\begin{frame}[t, plain]
  \plumheader{The Four Cases, Drawn}

  \begin{center}
  \begin{tikzpicture}[x=0.5cm, y=0.22cm]
    \draw[->, thick] (-2.7,0) -- (2.8,0);
    \draw[->, thick] (0,-5.0) -- (0,5.2);
    \draw[very thick, plum] plot [smooth] coordinates
      {(-2.2,4.07) (-2,0) (-1.7,-3.21) (-1.4,-4.0) (-1,-3) (-0.5,-0.94)
       (0,0) (0.5,-0.94) (1,-3) (1.4,-4.0) (1.7,-3.21) (2,0) (2.2,4.07)};
    \node[below, align=center, text width=3.2cm] at (0,-5.8)
      {\small \textbf{Case 1}\\ even, \(a_n > 0\)};
  \end{tikzpicture}
  \hspace{0.9cm}
  \begin{tikzpicture}[x=0.5cm, y=0.22cm]
    \draw[->, thick] (-2.7,0) -- (2.8,0);
    \draw[->, thick] (0,-5.0) -- (0,5.2);
    \draw[very thick, plum] plot [smooth] coordinates
      {(-2.2,-4.07) (-2,0) (-1.7,3.21) (-1.4,4.0) (-1,3) (-0.5,0.94)
       (0,0) (0.5,0.94) (1,3) (1.4,4.0) (1.7,3.21) (2,0) (2.2,-4.07)};
    \node[below, align=center, text width=3.2cm] at (0,-5.8)
      {\small \textbf{Case 2}\\ even, \(a_n < 0\)};
  \end{tikzpicture}
  \hspace{0.9cm}
  \begin{tikzpicture}[x=0.5cm, y=0.22cm]
    \draw[->, thick] (-2.7,0) -- (2.8,0);
    \draw[->, thick] (0,-5.0) -- (0,5.2);
    \draw[very thick, plum] plot [smooth] coordinates
      {(-2.2,-4.05) (-2,-2) (-1.8,-0.43) (-1.5,1.13) (-1,2) (-0.5,1.38)
       (0,0) (0.5,-1.38) (1,-2) (1.5,-1.13) (1.8,0.43) (2,2) (2.2,4.05)};
    \node[below, align=center, text width=3.2cm] at (0,-5.8)
      {\small \textbf{Case 3}\\ odd, \(a_n > 0\)};
  \end{tikzpicture}
  \hspace{0.9cm}
  \begin{tikzpicture}[x=0.5cm, y=0.22cm]
    \draw[->, thick] (-2.7,0) -- (2.8,0);
    \draw[->, thick] (0,-5.0) -- (0,5.2);
    \draw[very thick, plum] plot [smooth] coordinates
      {(-2.2,4.05) (-2,2) (-1.8,0.43) (-1.5,-1.13) (-1,-2) (-0.5,-1.38)
       (0,0) (0.5,1.38) (1,2) (1.5,1.13) (1.8,-0.43) (2,-2) (2.2,-4.05)};
    \node[below, align=center, text width=3.2cm] at (0,-5.8)
      {\small \textbf{Case 4}\\ odd, \(a_n < 0\)};
  \end{tikzpicture}
  \end{center}

  \bigskip
  \textit{An odd-degree polynomial has one end up and one end down, so its graph
  must cross the \(x\)-axis --- the same conclusion Lesson 2.4 reached by
  counting conjugate pairs.}
\end{frame}

% SECTION 4 — READING IT OFF ANY FORM
\begin{frame}[t, plain]
  \plumheader{4. Reading End Behavior Off Any Form}

  \begin{columns}[T]
    \column{0.48\textwidth}
      \begin{block}{Standard form hands it to you}
        \(p(x) = 9 + 2x^4 - x^7\)\\[4pt]
        Read the \textbf{exponents}, not the order. Leading term \(-x^7\):
        odd degree, negative \(a_n\).
        \[ \lim_{x \to \infty} p(x) = -\infty \]
        \[ \lim_{x \to -\infty} p(x) = \infty \]
      \end{block}

    \column{0.48\textwidth}
      \begin{block}{Factored form hides it --- never expand}
        \(q(x) = 2(x-3)^2(x+1)(x^2+4)\)\\[4pt]
        Degree: \(2 + 1 + 2 = 5\). Leading term:
        \(2 \cdot x^2 \cdot x \cdot x^2 = 2x^5\).
        \[ \lim_{x \to \infty} q(x) = \infty \]
        \[ \lim_{x \to -\infty} q(x) = -\infty \]
      \end{block}
  \end{columns}

  \bigskip
  \textit{Two forms, two jobs: factored form gave you the \textbf{zeros} in 2.4;
  standard form gives you the \textbf{end behavior} today. Lesson 2.6 shows they
  are the same polynomial.}
\end{frame}

% CLASS PRACTICE
\begin{frame}[t, plain]
  \plumheader{Class Practice}
  \sectionlabel[redacc]{Try It}

  \begin{enumerate}\setlength{\itemsep}{10pt}
    \item Write both limit statements for \(p(x) = -x^6 + 4x^3 - 11\).
    \item Write both limit statements for \(r(x) = -5(x+2)^3(x-1)^2\) --- do not
          expand.
    \item A polynomial has \(\lim_{x \to \infty} p(x) = \infty\) and
          \(\lim_{x \to -\infty} p(x) = -\infty\). Even or odd degree? Positive
          or negative \(a_n\)?
    \item For \(x^4 - 500x^3\), past what input is the leading term the larger
          of the two in size?
    \item Explain in one sentence why \(x^5 - 900x^4 + 12\) and \(x^5 + 7x\)
          have identical end behavior.
  \end{enumerate}
\end{frame}

\end{document}
